\subsection{半单模的张量积}

命题 4. —— 对 \(i \in \{1,2\}\)，设 \(\mathrm{A}_i\) 是一个 K-代数，\(\mathrm{M}_i\) 是一个半单 \(\mathrm{A}_i\)-模，\(\mathrm{Z}_i\) 是 \(\mathrm{M}_i\) 的交换子的中心。令 \(\mathrm{A} = \mathrm{A}_1 \otimes \mathrm{A}_2\)，\(\mathrm{M} = \mathrm{M}_1 \otimes \mathrm{M}_2\)，\(\mathrm{Z} = \mathrm{Z}_1 \otimes \mathrm{Z}_2\)。我们有 \(\Re_{\mathrm{A}}(\mathrm{M}) = \Re(\mathrm{Z})\mathrm{M}\)。若代数 \(\mathrm{Z}_1\) 与 \(\mathrm{Z}_2\) 之一在 \(\mathrm{K}\) 上可分，特别地若体 \(\mathrm{K}\) 是完备的，则 A-模 \(\mathrm{M}\) 无根。

对 \(i \in \{1,2\}\)，设 \(S_i\) 是一个单 \(A_i\)-模，\(D_i\) 是它的交换子，\(I(i)\) 是一个集合。我们先处理 \(M_i\) 是 \(A_i\)-模 \(S_i^{(I(i))}\) 的情形。我们把它的交换子的中心 \(Z_i\) 与 \(D_i\) 的中心视为同一。令 \(D = D_1 \otimes D_2\)。

我们有 \(\Re (\mathrm{D}) = \Re (\mathrm{Z})\mathrm{D}\)（VIII, 第 217 页，命题 3）与 \(\Re_{\mathrm{A}}(\mathrm{S}_1\otimes \mathrm{S}_2) =\) \(\Re (\mathrm{D})(\mathrm{S}_1\otimes \mathrm{S}_2)\)（VIII, 第 215 页，推论 2），因而 \(\Re_{\mathrm{A}}(\mathrm{S}_1\otimes \mathrm{S}_2) = \Re (\mathrm{Z})(\mathrm{S}_1\otimes \mathrm{S}_2)\)。A-模 M 是一族同构于 \(\mathrm{S}_1\otimes \mathrm{S}_2\) 的 A-模的直和，而一族模的直和的根是诸根的直和（VIII, 第 152 页，推论 2）。因此我们有 \(\Re_{\mathrm{A}}(\mathrm{M}) = \Re (\mathrm{Z})\mathrm{M}\)。

现在考虑一般情形。对 \(i \in \{1,2\}\)，把 \(A_{i}\)-模 \(M_{i}\) 的支集记作 \(S_{M_{i}}\)。对 \(\lambda \in S_{M_{i}}\)，把 \(M_{i}\) 的 \(\lambda\) 型同型分量记作 \(M_{i;\lambda}\)，并把它的交换子的中心记作 \(Z_{i;\lambda}\)。我们把环 \(Z_{i}\) 与诸环 \(Z_{i;\lambda}\)（\(\lambda \in S_{M_{i}}\)）的积视为同一。设 \(\lambda \in S_{M_{1}}\) 与 \(\mu \in S_{M_{2}}\)。用 \(i_{\lambda}: Z_{1;\lambda} \to Z_{1}\) 表示唯一的 K-线性映射，使得 \(pr_{\lambda} \circ i_{\lambda}\) 是 \(Z_{1;\lambda}\) 上的恒同映射，且 \(pr_{\lambda'} \circ i_{\lambda}\) 是零映射（对 \(\lambda' \in S_{M_{1}} = \{\lambda\}\)）。同样地定义 \(i_{\mu}: Z_{2;\mu} \to Z_{2}\)。令 \(Z_{\lambda,\mu} = Z_{1;\lambda} \otimes Z_{2;\mu}\) 与 \(i_{\lambda,\mu} = i_{\lambda} \otimes i_{\mu}\)，并把映射 \(pr_{\lambda} \otimes pr_{\mu}\)（从 Z 到 \(Z_{\lambda,\mu}\)）记作 \(\pi_{\lambda,\mu}\)。映射 \(\pi_{\lambda,\mu}\) 是一个满射环同态；因此我们有 \(\pi_{\lambda,\mu}(\Re(Z)) \subset \Re(Z_{\lambda,\mu})\)（VIII, 第 155 页，命题 5, b)）。我们来证明反包含。设 z 是 \(\Re(Z_{\lambda,\mu})\) 的一个元素；由于 \(Z_{1;\lambda}\) 与 \(Z_{2;\mu}\) 都是体，z 是幂零的（VIII, 第 215 页，定理 3）。我们有 \(i_{\lambda,\mu}(xy) = i_{\lambda,\mu}(x) \cdot i_{\lambda,\mu}(y)\)（对 \(Z_{\lambda,\mu}\) 中的 x、y）；因此 \(i_{\lambda,\mu}(z)\) 是幂零的，从而属于 \(\Re(Z)\)。由于 \(\pi_{\lambda,\mu} \circ i_{\lambda,\mu}\) 是 \(Z_{\lambda,\mu}\) 上的恒同映射，元素 z 属于 \(\pi_{\lambda,\mu}(\Re(Z))\)。我们就这样证明了等式 \(\pi_{\lambda,\mu}(\Re(Z)) = \Re(Z_{\lambda,\mu})\)。

令 \(M_{\lambda,\mu}=M_{1;\lambda}\otimes M_{2;\mu}\)；这是 M 的一个在 Z 下稳定的子模。对 \(z\in Z\) 与 \(m\in M_{\lambda,\mu}\)，我们有 \(zm=\pi_{\lambda,\mu}(z)m\)。因此，\(\Re(Z)M_{\lambda,\mu}\) 等于 \(\Re(Z_{\lambda,\mu})M_{\lambda,\mu}\)，从而由同型情形知等于 \(\Re_{A}(M_{\lambda,\mu})\)。由于直和的根是诸根的直和（VIII, 第 152 页，推论 2），且 M 是诸子模 \(M_{\lambda,\mu}\)（\((\lambda,\mu)\in\mathcal{S}_{\mathrm{M}_{1}}\times\mathcal{S}_{\mathrm{M}_{2}}\)）的直和，等式 \(\Re_{A}(M)=\Re(Z)M\) 得证。最后的断言则由 VIII, 第 217 页，命题 3 得出。

引理 2. —— 设 \(A_{1}\) 与 \(A_{2}\) 是交换域 K 上的代数。设 \(M_{1}\) 是一个在 K 上有限维的 \(A_{1}\)-模，\(M_{2}\) 是一个有限长度的 \(A_{2}\)-模。则 \(A_{1} \otimes A_{2}\)-模 \(M_{1} \otimes M_{2}\) 具有有限长度。

令 \(M = M_{1} \otimes M_{2}\)。设 \((e_{1}, \ldots, e_{n})\) 是 \(M_{1}\) 在体 K 上的一个基。映射 \((x_{1}, \ldots, x_{n}) \mapsto \sum_{i=1}^{n} e_{i} \otimes x_{i}\) 是从 \(A_{2}\)-模 \(M_{2}^{n}\) 到 \(A_{2}\)-模 M 的同构。由于 \(M_{2}\) 是有限长度的 \(A_{2}\)-模，M 亦然。此外，M 的每个 A-子模都是 \(A_{2}\)-子模；因此 M 是有限长度的 A-模。

命题 5. —— 设 \(A_{1}\) 与 \(A_{2}\) 是交换域 K 上的代数。设 \(M_{1}\) 是一个在 K 上有限维的半单 \(A_{1}\)-模，\(M_{2}\) 是一个半单 \(A_{2}\)-模。对 \(i \in \{1,2\}\)，把 \(A_{i}\)-模 \(M_{i}\) 的交换子记作 \(D_{i}\)，并把 \(D_{i}\) 的中心记作 \(Z_{i}\)。令 \(A = A_{1} \otimes A_{2}\)，\(M = M_{1} \otimes M_{2}\)，\(D = D_{1} \otimes D_{2}\)，\(Z = Z_{1} \otimes Z_{2}\)。

\begin{enumerate}
\def\labelenumi{\alph{enumi})}
\item
  A-模 M 的交换子可与 D 视为同一，它的中心是 Z。若 \(A_{2}\)-模 \(M_{2}\) 具有有限长度，则 A-模 M 具有有限长度，环 D 是右且左阿廷的，环 Z 是阿廷的。
\item
  下列性质等价：
\end{enumerate}

\begin{enumerate}
\def\labelenumi{(\roman{enumi})}
\item
  A-模 M 是半单的。
\item
  环 \(Z\) 同构于一族交换域的积。
\item
  环 Z 是既约的。
\end{enumerate}

\begin{enumerate}
\def\labelenumi{\alph{enumi})}
\setcounter{enumi}{2}
\tightlist
\item
  下列性质等价：
\end{enumerate}

\begin{enumerate}
\def\labelenumi{(\roman{enumi})}
\item
  A-模 M 是同型的且不约化为 0。
\item
  环 Z 是一个体。
\item
  环 \(\mathbf{Z}\) 是一个整环。
\end{enumerate}

由假设，\(M_{1}\) 在 K 上有限维。因此 M 的交换子可与 D 视为同一（VIII, 第 213 页，引理 1），且它的中心是 Z（III, §4, No.~4, 第 468 页，推论）。假设 \(A_{2}\)-模 \(M_{2}\) 具有有限长度。由引理 2，A-模 M 具有有限长度。由于 \(M_{2}\) 是半单且有限生成的，环 \(D_{2}\) 是半单的（VIII, 第 139 页，命题 6），且它的中心 \(Z_{2}\) 是有限族交换域的积。于是，\(D_{2}\)-模 \((\mathrm{D}_{2})_{s}\) 与 \(Z_{2}\)-模 \((\mathrm{Z}_{2})_{s}\) 都具有有限长度。此外，由于 \(M_{1}\) 在 K 上有限维，\(D_{1}\) 与 \(Z_{1}\) 亦然。由引理 2，模 \((\mathrm{D}_{1} \otimes \mathrm{D}_{2})_{s}\) 具有有限长度；因此环 \(D_{1} \otimes D_{2}\) 是左阿廷的。我们同样可证环 \(D_{1} \otimes D_{2}\) 是右阿廷的，且环 \(Z_{1} \otimes Z_{2}\) 是阿廷的。断言 a) 得证。

我们来证明 b)。半单模的交换子的中心同构于一族交换域的积（VIII, 第 87 页，命题 8, a)）；这就证明了 (i) 蕴含 (ii)。蕴含关系 (ii) \(\Rightarrow\) (iii) 是明显的。

假设环 Z 是既约的。则我们有 \(\Re(Z)=0\)（VIII, 第 215 页，定理 3）与 \(\Re_{\mathrm{A}}(\mathrm{M})=0\)（VIII, 第 218 页，命题 4）。由于 \(A_{2}\)-模 \(M_{2}\) 是半单的，存在一个族 \((\mathrm{S}_{i})_{i\in\mathrm{I}}\)（由单 \(A_{2}\)-模组成）以及一个从 \(M_{2}\) 到 \(\bigoplus S_{i}\) 的同构。因此，A-模 M 同构于 \(\bigoplus M_{1}\otimes S_{i}\)。于是对每个 \(i\in I\)，A-模 \(M_{1}\otimes S_{i}\) 无根。由 a)，它具有有限长度；故它是半单的（VIII, 第 153 页，命题 3, b)）。于是 A-模 M 是一族半单模的直和，从而是半单的。这就证明了 (iii) 蕴含 (i)，并完成了 b) 的证明。

一个 A-模是同型的且非零的，当且仅当它是半单的且其交换子的中心是一个体（VIII, 第 87 页，命题 8, b)）。于是 c) 由 b) 得出。

推论. —— 若 \(Z_{1}\) 或 \(Z_{2}\) 是体 K 上的可分代数（例如当 K 完备时即属此种情形），则 A-模 \(M_{1} \otimes M_{2}\) 是半单的。

环 \(Z_{1}\) 与 \(Z_{2}\) 同构于诸体的积，因而是既约环。特别地，若 K 是完备的，则它们是 K 上的可分代数（V, §15, No.~5, 第 125 页，定理 3）。由 V, §15, No.~2, 第 120 页，命题 5，可分代数与既约代数的张量积是既约的。于是 Z 是既约环，推论由命题 5, b) 得出。

